\hypertarget{filesystem}{%
\chapter{4. Filesystem}\label{filesystem}}

The filesystem is the widest surface the workload touches. The user's
keys, tokens, browser state, shell history, configuration, and work in
progress all live on it, and under the default-uid model a workload
running as the user reaches every byte. A path not meant for the
workload is only a path it has not read yet.

Of the two limbs, the filesystem yields almost entirely to the first.
Most of what a workload must be kept from is data it has no reason to
touch, and the surest way to keep it from that data is for the data not
to be in its world at all (\texttt{do-less}). So the filesystem is
\texttt{construction-by-absence} worked to its edges: not a filter that
checks each read against a list of the forbidden, but a world built so
the forbidden is not there to read.

\hypertarget{the-empty-world}{%
\section{4.1 The empty world}\label{the-empty-world}}

The framework builds the workload an empty world and places into it
exactly what the policy grants. It does not hand over the user's home
and then block the parts that should stay hidden; what the policy does
not grant is not forbidden, it is missing.

The difference is not cosmetic, and it shows up most clearly against
reconnaissance. A guard over the user's home leaks the home's shape: a
workload that asks for a secret and is refused has learned the secret is
there to be refused. An agent mapping the environment for credentials
and history (\texttt{T1.1}), or a post-install script reaching for a key
ring (\texttt{T1.2}), can read the topology of the user's life off the
pattern of what it is allowed to open. Refusal is an answer. Absence is
silence. In the workload's world the sensitive directories do not exist,
so the agent meets no closed door: it meets nothing, and learns nothing,
because there is nothing to learn.

The boundary also costs nothing to hold. No monitor reads over the
workload's shoulder; nothing is checked on each open. What the workload
can reach was settled when its world was built, before it ran, and the
operating system holds it there. Reachability is the shape of the world,
not a decision remade on every access.

\hypertarget{the-contract}{%
\section{4.2 The contract}\label{the-contract}}

The contract is default-deny, and the denial is structural rather than
enforced: an empty policy yields an empty world. That world is still
workable. The system's own shared files, the libraries and programs
every process needs in order to run, are present and read-only, because
they carry nothing of the user's; what is omitted is everything that
does. A grant then adds back exactly one of two things: a path the
workload may read and traverse, or a path it may write.

The workload acts under the operator's own identity. It has to (a
granted path is the operator's, and only the operator's identity opens
it), but it acts with none of the authority that identity ordinarily
carries. Its reach is reduced to the paths the policy names, and the
rest of what the operator could touch is not fenced off; it is absent.
This is the split the whole design turns on (\texttt{split-the-uid}):
the identity kept, the authority pared back to the grant. Even the
operator's name is pared away: the workload sees a generic, confined
identity, not whose machine it is running on, so the question of whose
data this is goes unanswered along with the data.

Writing follows the same rule as reading: it lands where the policy
grants it and nowhere else. The world is writable, but only the paths
granted to persist outlive the workload; a file written anywhere else is
gone when the workload ends. A confined workload cannot litter the host
with leavings, and it cannot leave anything behind outside the
directories it was told to work in. Its scratch space is its own: the
host's shared temporary area, with other sessions' intermediate files
and other users' leavings, is absent, replaced by a private space
destroyed along with the workload.

One thing a filesystem grant is not is a grant of hardware. A device is
not a path, whatever it looks like in a directory listing; handing a
workload raw memory or a disk is handing it the machine, not a file, and
the filesystem policy will not do it. Hardware stays below this boundary
and is answered on its own terms.

\hypertarget{absence-at-finer-grain}{%
\section{4.3 Absence at finer grain}\label{absence-at-finer-grain}}

Some files cannot be removed, because the workload needs them. A package
tool needs the configuration that tells it where to find its registries;
a version-control client needs the user's global settings. The trouble
is that those same files often carry a secret a line or two from the
setting: an auth token in the registry configuration, a credential
helper in the version-control settings (\texttt{T1.1}). Removing the
file breaks the tool. Handing it over whole leaks the secret.

So absence goes finer. Where it cannot take the file, it takes the
secret. A configuration the workload needs is placed in its world with
the secret-bearing lines gone: the settings that make the tool work
remain, and the credentials are simply not in the copy the workload
reads. A secret that lives as its own file inside an otherwise-granted
directory (an environment file, a key beside the code that uses it) is
present but empty: there, so a program that expects it does not fall
over, but holding nothing. An agent walking the project tree for an
environment file finds it, opens it, and reads no password, because the
password was never placed in the world it is reading.

This is the same move as the empty world, one level down. Reconnaissance
inside a granted directory comes back as empty-handed as reconnaissance
outside one.

\hypertarget{the-floor}{%
\section{4.4 The floor}\label{the-floor}}

All of this assumes the policy is written well. The design does not
assume it.

Beneath every policy sits a floor the policy cannot lower. A set of
paths is categorically forbidden: the SSH and signing keys, the cloud
credentials, the password store, the browser profiles, the system's own
secret files, and they are forbidden before any grant is considered, not
subtracted after. A policy that, through haste or a copied template,
grants read access to the whole of the user's home still does not expose
them, because the framework will not build a world in which they are
reachable at all. Where a grant and the floor disagree, the floor wins.

This is fail-safe defaults {[}SaltzerSchroeder{]} taken at its word: the
safe outcome does not depend on the operator getting the policy right.
The most sensitive data on the machine is absent by default and stays
absent in spite of the policy, not because of it.

The floor is also where the boundary is held a second time. The empty
world removes the forbidden by construction, but construction can be
probed: a path assembled through a traversal trick, a symbolic link
aimed past the grant. So the same boundary is enforced twice over: the
world is built without the forbidden in it, and the operating system
independently refuses any read that resolves outside the granted paths.
A workload that finds a seam in the first runs straight into the second.

\hypertarget{the-trust-record}{%
\section{4.5 The trust record}\label{the-trust-record}}

There is one write the empty world does not contain on its own: a write
to a file the operator will later run.

A workload with write access to a project directory can change more than
data. It can edit a build script, a task definition, a source-control
hook: files that do nothing inside the kennel but execute with the
operator's full, unconfined authority the next time the operator builds
the project or commits to it. The confined workload plants the trap; the
unconfined operator springs it (\texttt{T2.8}). The filesystem boundary
does not catch this on its own, because the file sits inside a directory
the workload was legitimately granted.

The design answers it by recording, before the workload runs, a trusted
baseline of every file in the workspace that can trigger execution: the
scripts, the hooks, the task definitions. That baseline is the reference
the operator will later judge the workload's changes against. For it to
mean anything, the workload must not be able to forge it, and the
workload has write access to the very directory the baseline describes.
So absence is turned on the record itself. The baseline is absent from
the workload's world: the workload can rewrite a build script, but it
cannot see, read, or reconstruct the record of what that script was
trusted to be. It can change the trigger. It cannot change the pin.

When the workload finishes, the operator reviews. Host-side tooling,
which holds the real baseline the workload never saw, compares the
workspace's triggers against it and shows what moved. The operator
accepts the changes, or has the triggers restored to their recorded
state, erasing a planted payload before it can run. The workload's
writes are held to data, and the path from a confined write to
unconfined execution is cut.

Across all of it the filesystem is the first limb and little else. There
is no per-read broker and no transaction on each open, because none is
needed: a resource that is absent needs no mediating, and the discipline
of the chapter is to make as much absent as it can: whole subtrees, then
single files, then the secret lines within a file, down to the record
that guards the workload's own writes. What the workload cannot reach it
cannot probe, and what it cannot see it cannot forge. The probe comes
back empty because there is nothing to find.
